\section*{Đề thi thử THPTQG Sở GD\&ĐT An Giang 25--26, Lần 1}

\caulc
\Opensolutionfile{ans}[ans/ans-56-26-2-AnGiang-SoGD-L1-LC]
\begin{ex}%[Thi thử SGD An Giang 25-26,L1]%[Lê Quốc Hiệp, DeTHPT2025-2026 ]%[2D4N1-4]
	Họ nguyên hàm của hàm số $f(x)=5\cdot 2^x$ là
	\choice
	{\True $\dfrac{5\cdot 2^x}{\ln 2}+C$}
	{$5\cdot 2^x\ln 2+C$}
	{$\dfrac{2^x}{\ln 2}+C$}
	{$\dfrac{5\cdot 2^{x+1}}{x+1}+C$}
	\loigiai{
		Ta có $\displaystyle\int f(x)\,\mathrm{d}x=\displaystyle\int 5\cdot 2^x\,\mathrm{d}x=5\cdot\dfrac{2^x}{\ln 2}+C=\dfrac{5\cdot 2^x}{\ln 2}+C$.
	}
\end{ex}
\begin{ex}%[Thi thử SGD An Giang 25-26,L1]%[Lê Quốc Hiệp, DeTHPT2025-2026 ]%[2H5N2-2]
	Trong không gian với hệ tọa độ $Oxyz$, cho đường thẳng $(d)\colon\dfrac{x-1}{2}=\dfrac{y+3}{-4}=\dfrac{z-5}{3}$. Vectơ nào sau đây là một vectơ chỉ phương của đường thẳng $(d)$?
	\choice
	{$\vec{u}_1=(1;-3;5)$}
	{$\vec{u}_2=(2;4;3)$}
	{\True $\vec{u}_3=(2;-4;3)$}
	{$\vec{u}_4=(-2;-4;-3)$}
	\loigiai{
		Đường thẳng $(d):\dfrac{x-1}{2}=\dfrac{y+3}{-4}=\dfrac{z-5}{3}$ có một vectơ chỉ phương là $\vec{u}=(2;-4;3)$.
		Đối chiếu với các phương án, ta chọn $\vec{u}_3=(2;-4;3)$.
	}
\end{ex}
\begin{ex}%[Thi thử SGD An Giang 25-26,L1]%[Lê Quốc Hiệp, DeTHPT2025-2026 ]%[2H2H1-2]
	\immini
	{Cho hình chóp $S.ABCD$ có đáy $ABCD$ là hình bình hành (xem hình bên). Gọi $O$ là giao điểm của hai đường chéo. Khẳng định nào sau đây là đúng?
	\choice
	{$\overrightarrow{SA}+\overrightarrow{SB}=\overrightarrow{SC}+\overrightarrow{SD}$}
	{$\overrightarrow{AB}+\overrightarrow{BC}=\overrightarrow{DB}$}
	{\True $\overrightarrow{SA}+\overrightarrow{SC}=\overrightarrow{SB}+\overrightarrow{SD}$}
	{$\overrightarrow{SA}+\overrightarrow{SB}+\overrightarrow{SC}+\overrightarrow{SD}=\overrightarrow{0}$}}
	{\begin{tikzpicture}[line cap=round,line join=round,font=\footnotesize,>=stealth,scale=1]
			\fill (0,0) coordinate [label=above left:$A$] (A) circle(1pt)
			(0:3) coordinate [label=below right:$B$] (B) circle(1pt)
			(90:2) coordinate [label=above right:$S$] (S) circle(1pt)
			(-40:2) coordinate [label=below right:$C$] (C) circle(1pt)
			($(A)+(C)-(B)$) coordinate [label=below:$D$] (D) circle(1pt)
			($0.5*(C)$) coordinate [label=below:$O$] (O) circle(1pt);
			\draw (S)--(B)--(C) (D)--(S)--(C)--(D);
			\draw[dashed] (D)--(A)--(S) (C)--(A)--(B)--(D);
	\end{tikzpicture}}
	\loigiai{
		Vì $ABCD$ là hình bình hành tâm $O$ nên $O$ là trung điểm của hai đường chéo $AC$ và $BD$.
		Xét tam giác $SAC$ có trung tuyến $SO$, ta có $\overrightarrow{SA}+\overrightarrow{SC}=2\overrightarrow{SO}$.
		Xét tam giác $SBD$ có trung tuyến $SO$, ta có $\overrightarrow{SB}+\overrightarrow{SD}=2\overrightarrow{SO}$.
		Suy ra $\overrightarrow{SA}+\overrightarrow{SC}=\overrightarrow{SB}+\overrightarrow{SD}$.
	}
\end{ex}
\begin{ex}%[Thi thử SGD An Giang 25-26,L1]%[Lê Quốc Hiệp, DeTHPT2025-2026 ]%[2H2N2-3]
	Trong không gian với hệ trục tọa độ $Oxyz$, cho vectơ $\vec{u}=-3\vec{i}+5\vec{j}+2\vec{k}$. Tọa độ của vectơ $\vec{u}$ là
	\choice
	{$(3;5;2)$}
	{\True $(-3;5;2)$}
	{$(-3;5;-2)$}
	{$(5;-3;2)$}
	\loigiai{
		Ta có $\vec{u}=-3\vec{i}+5\vec{j}+2\vec{k} \Rightarrow \vec{u}=(-3;5;2)$.
	}
\end{ex}
\begin{ex}%[Thi thử SGD An Giang 25-26,L1]%[Lê Quốc Hiệp, DeTHPT2025-2026 ]%[1D1N5-3]
	Tập nghiệm của phương trình $\cos x=0$ là
	\choice
	{$S=\{k2\pi \mid k\in\mathbb{Z}\}$}
	{$S=\{k\pi \mid k\in\mathbb{Z}\}$}
	{$S=\left\{\dfrac{\pi}{2}+k2\pi \mid k\in\mathbb{Z}\right\}$}
	{\True $S=\left\{\dfrac{\pi}{2}+k\pi \mid k\in\mathbb{Z}\right\}$}
	\loigiai{
		Ta có $\cos x=0 \Leftrightarrow x=\dfrac{\pi}{2}+k\pi$ ($k\in\mathbb{Z}$).
		Vậy tập nghiệm của phương trình là $S=\left\{\dfrac{\pi}{2}+k\pi \mid k\in\mathbb{Z}\right\}$.
	}
\end{ex}
\begin{ex}%[Thi thử SGD An Giang 25-26,L1]%[Lê Quốc Hiệp, DeTHPT2025-2026 ]%[1D2N2-4]
	Cho cấp số cộng $(u_n)$ có $u_1=3$ và công sai $d=-2$. Giá trị của $u_7$ bằng
	\choice
	{$-11$}
	{\True $-9$}
	{$15$}
	{$-14$}
	\loigiai{
		Áp dụng công thức số hạng tổng quát của cấp số cộng, ta có:
		\[u_7=u_1+6d=3+6\cdot(-2)=-9.\]
	}
\end{ex}
\begin{ex}%[Thi thử SGD An Giang 25-26,L1]%[Lê Quốc Hiệp, DeTHPT2025-2026 ]%[1D6N4-2]
	Nghiệm của phương trình $\log_2(3x-1)=3$ là
	\choice
	{$x=2$}
	{\True $x=3$}
	{$x=\dfrac{10}{3}$}
	{$x=4$}
	\loigiai{
		Điều kiện: $3x-1>0 \Leftrightarrow x>\dfrac{1}{3}$.\\
		Ta có $\log_2(3x-1)=3 \Leftrightarrow 3x-1=2^3 \Leftrightarrow 3x-1=8 \Leftrightarrow 3x=9 \Leftrightarrow x=3$ (thỏa mãn điều kiện).\\
		Vậy phương trình có nghiệm là $x=3$.
	}
\end{ex}
\begin{ex}%[Thi thử SGD An Giang 25-26,L1]%[Lê Quốc Hiệp, DeTHPT2025-2026 ]%[2D1H1-1]
	Cho hàm số $y=f(x)$ có đạo hàm $f'(x)$ liên tục trên $\mathbb{R}$. Biết rằng $f'(x)=(x^2+1)(x-2)$ với mọi $x\in\mathbb{R}$. Hàm số $y=f(x)$ đồng biến trên khoảng nào dưới đây?
	\choice
	{$(0;2)$}
	{$(-\infty;-1)$}
	{\True $(2;+\infty)$}
	{$(-1;2)$}
	\loigiai{
		Ta có $f'(x)=0\Leftrightarrow(x^2+1)(x-2)=0\Leftrightarrow x=2$ (do $x^2+1>0$ với mọi $x\in\mathbb{R}$).\\
		Hàm số $y=f(x)$ đồng biến khi $f'(x)>0\Leftrightarrow x-2>0\Leftrightarrow x>2$.\\
		Vậy hàm số đã cho đồng biến trên khoảng $(2;+\infty)$.
	}
\end{ex}
\begin{ex}%[Thi thử SGD An Giang 25-26,L1]%[Lê Quốc Hiệp, DeTHPT2025-2026 ]%[1D5H2-2]
	Một người thống kê thời lượng (đơn vị: giây) thực hiện các cuộc gọi điện thoại của mình trong một tuần và lập bảng tần số ghép nhóm như sau:
	\begin{center}
		\begin{tabular}{|c|c|c|c|c|c|c|}
			\hline
			Nhóm & $[0;40)$ & $[40;80)$ & $[80;120)$ & $[120;160)$ & $[160;200)$ & $[200;240)$ \\
			\hline
			Tần số & $8$ & $10$ & $12$ & $6$ & $3$ & $1$ \\
			\hline
		\end{tabular}
	\end{center}
	Trung vị $M_e$ (đơn vị: giây) của mẫu số liệu ghép nhóm trên bằng
	\choice
	{$84$}
	{\True $\dfrac{260}{3}$}
	{$80$}
	{$86{,}5$}
	\loigiai{
		Số phần tử của mẫu là $n=8+10+12+6+3+1=40$. Ta có
		\[\dfrac{n}{2}=\dfrac{40}{2}=20.\]
		\begin{center}
			\begin{tabular}{|c|c|c|c|c|c|c|}
				\hline
				Nhóm & $[0;40)$ & $[40;80)$ & $[80;120)$ & $[120;160)$ & $[160;200)$ & $[200;240)$ \\
				\hline
				Tần số & $8$ & $10$ & $12$ & $6$ & $3$ & $1$ \\
				\hline
				Tần số tích lũy & $8$ & $18$ & $30$ & $36$ & $39$ & $40$ \\
				\hline
			\end{tabular}
		\end{center}
		Gọi $cf_2$, $cf_3$ lần lượt là tần số tích luỹ của nhóm $2$ và nhóm $3$. Ta có $cf_2=18<20<cf_3=30$.\\
		Suy ra nhóm $3$ là nhóm đầu tiên có tần số tích luỹ lớn hơn hoặc bằng $20$.\\
		Do đó, nhóm chứa trung vị là nhóm $[80;120)$.\\
		Xét nhóm $3$ là nhóm $[80;120)$ có $r=80$, $d=40$, $n_3=12$ và nhóm $2$ là nhóm $[40;80)$ có $cf_2=18$.\\
		Áp dụng công thức, ta có trung vị của mẫu số liệu là
		\[M_e=80+\dfrac{20-18}{12}\cdot40=80+\dfrac{20}{3}=\dfrac{260}{3}.\]
	}
\end{ex}
\begin{ex}%[Thi thử SGD An Giang 25-26,L1]%[Lê Quốc Hiệp, DeTHPT2025-2026 ]%[2H5N2-2]
	Trong không gian $Oxyz$, cho hai điểm $M(2;1;-1)$ và $N(4;-1;0)$. Một vectơ chỉ phương của đường thẳng $MN$ là
	\choice
	{$\vec{u}=(6;0;-1)$}
	{\True $\vec{u}=(2;-2;1)$}
	{$\vec{u}=(2;2;1)$}
	{$\vec{u}=(-2;2;1)$}
	\loigiai{
		Ta có
		\[\overrightarrow{MN}=(4-2;-1-1;0-(-1))=(2;-2;1).\]
		Vậy một vectơ chỉ phương của đường thẳng $MN$ là $\vec{u}=\overrightarrow{MN}=(2;-2;1)$.
	}
\end{ex}
\begin{ex}%[Thi thử SGD An Giang 25-26,L1]%[Lê Quốc Hiệp, DeTHPT2025-2026 ]%[1H8H5-3]
	\immini{
		Cho hình lập phương $ABCD.A'B'C'D'$ có cạnh bằng $a$. Khoảng cách từ điểm $A$ đến mặt phẳng $(BDD'B')$ bằng
		\choice
		{$a\sqrt{2}$}
		{$a$}
		{\True $\dfrac{a\sqrt{2}}{2}$}
		{$\dfrac{a\sqrt{3}}{2}$}
	}{
		\begin{tikzpicture}[scale=1, font=\footnotesize, line join=round, line cap=round, >=stealth]
			\def\a{3}
			\coordinate (A) at (0,0);
			\coordinate (B) at (\a,0);
			\coordinate (D) at (0.6,0.8);
			\coordinate (C) at ($ (B) + (D) - (A) $);
			\coordinate (A') at ($ (A) + (0,\a) $);
			\coordinate (B') at ($ (B) + (0,\a) $);
			\coordinate (C') at ($ (C) + (0,\a) $);
			\coordinate (D') at ($ (D) + (0,\a) $);
			\fill[blue!10, opacity=0.8] (B) -- (D) -- (D') -- (B') -- cycle;
			\draw (A) node[below left]{$A$} -- (B) node[below right]{$B$} -- (C) node[right]{$C$} -- (C') node[right]{$C'$} -- (B') node[right]{$B'$} -- (A') node[left]{$A'$} -- cycle;
			\draw (A) -- (A');
			\draw (B) -- (B');
			\draw (C') -- (D') node[above left]{$D'$} -- (A');
			\draw (B') -- (D');
			\draw[dashed] (A) -- (D) node[above left]{$D$} -- (C);
			\draw[dashed] (D) -- (D');
			\draw[dashed] (B) -- (D);
			\draw[dashed] (A) -- (C);
		\end{tikzpicture}
	}
	\loigiai{
		\immini
		{Gọi $O$ là giao điểm của $AC$ và $BD$.
		Vì $ABCD$ là hình vuông nên $AO\perp BD$ (1).\\
		Mặt khác, $BB'\perp (ABCD)\Rightarrow BB'\perp AO$ (2).\\
		Từ (1) và (2) suy ra $AO\perp(BDD'B')$.\\
		Do đó, khoảng cách từ $A$ đến mặt phẳng $(BDD'B')$ chính là độ dài đoạn thẳng $AO$.\\
		Trong hình vuông $ABCD$ cạnh $a$, ta có $AC=a\sqrt{2}$.\\
		Vậy khoảng cách cần tìm là $AO=\dfrac{AC}{2}=\dfrac{a\sqrt{2}}{2}$.}
		{\begin{tikzpicture}[scale=1, font=\footnotesize, line join=round, line cap=round, >=stealth]
				\def\a{3}
				\coordinate (A) at (0,0);
				\coordinate (B) at (\a,0);
				\coordinate (D) at (0.7,0.8);
				\coordinate (C) at ($ (B) + (D) - (A) $);
				\coordinate (O) at ($ 0.5*(C)$);
				\coordinate (A') at ($ (A) + (0,\a) $);
				\coordinate (B') at ($ (B) + (0,\a) $);
				\coordinate (C') at ($ (C) + (0,\a) $);
				\coordinate (D') at ($ (D) + (0,\a) $);
				
				\fill[blue!10, opacity=0.8] (B) -- (D) -- (D') -- (B') -- cycle;
				
				\draw (A) node[below left]{$A$} -- (B) node[below right]{$B$} -- (C) node[right]{$C$} -- (C') node[right]{$C'$} -- (B') node[right]{$B'$} -- (A') node[left]{$A'$} -- cycle;
				\draw (A) -- (A');
				\draw (B) -- (B');
				\draw (C') -- (D') node[above left]{$D'$} -- (A');
				\draw (B') -- (D');
				\draw[dashed] (A) -- (D) node[above left]{$D$} -- (C);
				\draw[dashed] (D) -- (D');
				\draw[dashed] (B) -- (D);
				\draw[dashed] (A) -- (C);
				\draw (O) node[below=-0.5pt]{$O$};
				% Thêm chấm các điểm bằng vòng lặp
				\foreach \P in {A,B,C,D,A',B',C',D',O}
				\fill (\P) circle (1pt);
		\end{tikzpicture}}
	}
\end{ex}
\begin{ex}%[Thi thử SGD An Giang 25-26,L1]%[Lê Quốc Hiệp, DeTHPT2025-2026 ]%[2D4H3-3]
	Thể tích vật thể tròn xoay sinh ra khi quay hình phẳng giới hạn bởi đồ thị $y=x^2$, trục hoành và hai đường thẳng $x=0$, $x=2$ quanh trục hoành là
	\choice
	{\True $\dfrac{32\pi}{5}$}
	{$\dfrac{8\pi}{3}$}
	{$\dfrac{32}{5}$}
	{$\dfrac{16\pi}{5}$}
	\loigiai{
		Thể tích vật thể tròn xoay cần tìm là
		\[V=\pi\displaystyle\int\limits_0^2(x^2)^2\,\mathrm{d}x=\pi\displaystyle\int\limits_0^2x^4\,\mathrm{d}x=\left.\dfrac{\pi x^5}{5}\right|_0^2=\dfrac{32\pi}{5}.\]
	}
\end{ex}
\Closesolutionfile{ans}


\cauds
\Opensolutionfile{ans}[ans/ans-56-26-2-AnGiang-SoGD-L1-DS]
\begin{ex}%[Thi thử SGD An Giang 25-26,L1]%[Lê Quốc Hiệp, DeTHPT2025-2026 ]%[2D1H3-1]
	Cho hàm số $f(x)=x^3-12x+5$.
	\choiceTF
	{\True Hàm số đã cho có đạo hàm là $f'(x)=3x^2-12$}
	{Phương trình $f'(x)=0$ có tập nghiệm là $S=\{2\}$}
	{$f(2)=11$}
	{Giá trị lớn nhất của hàm số $f(x)$ trên đoạn $[-3;3]$ bằng $14$}
	\loigiai{
		\begin{itemchoice}
			\itemch Ta có $f'(x)=(x^3-12x+5)'=3x^2-12$.
			\itemch Ta có $f'(x)=0\Leftrightarrow3x^2-12=0\Leftrightarrow x=\pm2$. \\
			Vậy tập nghiệm của phương trình là $S=\{-2;2\}$.
			\itemch Ta có $f(2)=2^3-12\cdot2+5=-11$.
			\itemch Hàm số $f(x)$ liên tục trên đoạn $[-3;3]$.\\
			Ta có $f'(x)=0\Leftrightarrow x=\pm2\in[-3;3]$.\\
			Ta có
			\begin{eqnarray*}
				f(-3)&=&(-3)^3-12\cdot(-3)+5=14.\\
				f(-2)&=&(-2)^3-12\cdot(-2)+5=21.\\
				f(2)&=&-11.\\
				f(3)&=&3^3-12\cdot3+5=-4.
			\end{eqnarray*}
			Vậy giá trị lớn nhất của hàm số $f(x)$ trên đoạn $[-3;3]$ bằng $21$.
		\end{itemchoice}
	}
\end{ex}
\begin{ex}%[Thi thử SGD An Giang 25-26,L1]%[Lê Quốc Hiệp, DeTHPT2025-2026 ]%[2D4H1-6]
	Một quần thể vi khuẩn ban đầu có $1000$ con. Sau $2$ giờ, số lượng vi khuẩn tăng lên thành $4000$ con. Gọi $P(t)$ là số lượng vi khuẩn tại thời điểm $t$ (giờ). Biết rằng tốc độ tăng trưởng của quần thể tỉ lệ thuận với số lượng vi khuẩn hiện có, tức là $P'(t)=k\cdot P(t)$, với $k$ là hằng số khác $0$ và $P(t)>0$ với $t\ge0$.
	\choiceTF
	{\True Hằng số tốc độ tăng trưởng của quần thể vi khuẩn là $k=\ln 2$}
	{\True Sau $5$ giờ kể từ thời điểm ban đầu, số lượng vi khuẩn là $32000$ con}
	{Công thức tính số lượng vi khuẩn tại thời điểm $t$ là $P(t)=4000\mathrm{e}^{kt}$ với $t\ge0$}
	{Để số lượng vi khuẩn đạt $64000$ con, cần đúng $8$ giờ kể từ thời điểm ban đầu}
	\loigiai{
		Ta có
		\begin{eqnarray*}
			&&P'(t)=k\cdot P(t)\\
			&\Leftrightarrow&\dfrac{P'(t)}{P(t)}=k\\
			&\Leftrightarrow&\displaystyle\int\dfrac{P'(t)}{P(t)}\,\mathrm{d}t=\displaystyle\int k\,\mathrm{d}t\\
			&\Leftrightarrow&\ln P(t)=kt+C\\
			&\Leftrightarrow& P(t)=\mathrm{e}^{kt+C}=A\cdot \mathrm{e}^{kt}\text{ (với }A=\mathrm{e}^C>0\text{)}.
		\end{eqnarray*}
		Ban đầu ($t=0$) có $1000$ con nên $P(0)=1000\Rightarrow A\cdot \mathrm{e}^0=1000\Rightarrow A=1000$.\\
		Suy ra công thức tính số lượng vi khuẩn là $P(t)=1000\mathrm{e}^{kt}$.\\
		Sau $2$ giờ, số lượng vi khuẩn là $4000$ con nên
		\[P(2)=4000\Rightarrow1000\mathrm{e}^{2k}=4000\Leftrightarrow \mathrm{e}^{2k}=4\Leftrightarrow2k=\ln 4\Leftrightarrow k=\ln 2.\]
		Khi đó, $P(t)=1000\mathrm{e}^{t\ln 2}=1000\cdot2^t$.
		\begin{itemchoice}
			\itemch Hằng số tốc độ tăng trưởng $k=\ln 2$.
			\itemch Sau $5$ giờ, số lượng vi khuẩn là $P(5)=1000\cdot2^5=32000$ (con).
			\itemch Công thức tính số lượng vi khuẩn là $P(t)=1000\mathrm{e}^{kt}$.
			\itemch Để số lượng vi khuẩn đạt $64000$ con thì $P(t)=64000\Leftrightarrow1000\cdot2^t=64000\Leftrightarrow2^t=64\Leftrightarrow t=6$ (giờ).
		\end{itemchoice}
	}
\end{ex}
\begin{ex}%[Thi thử SGD An Giang 25-26,L1]%[Lê Quốc Hiệp, DeTHPT2025-2026 ]%[2D6H2-4]
	Một công ty điện tử có hai phân xưởng $I$ và $II$ cùng sản xuất một loại linh kiện. Biết rằng phân xưởng $I$ sản xuất $75\%$ tổng số lượng linh kiện của công ty, phân xưởng $II$ sản xuất $25\%$ còn lại. Tỉ lệ linh kiện bị lỗi của phân xưởng $I$ và phân xưởng $II$ lần lượt là $4\%$ và $2\%$. Chọn ngẫu nhiên một linh kiện trong kho của công ty để kiểm tra.
	\choiceTF
	{Nếu linh kiện được chọn do phân xưởng $I$ sản xuất thì xác suất để linh kiện đó không bị lỗi là $0{,}94$}
	{\True Xác suất để linh kiện được chọn bị lỗi và do phân xưởng $I$ sản xuất là $0{,}03$}
	{\True Xác suất để linh kiện được chọn bị lỗi là $0{,}035$}
	{Nếu linh kiện được chọn bị lỗi thì xác suất để linh kiện đó do phân xưởng $I$ sản xuất là $\dfrac{4}{7}$}
	\loigiai{
		Gọi $A_1$ là biến cố ``Linh kiện được chọn do phân xưởng $I$ sản xuất''. Ta có $\mathrm{P}(A_1)=0{,}75$.\\
		Gọi $A_2$ là biến cố ``Linh kiện được chọn do phân xưởng $II$ sản xuất''. Ta có $\mathrm{P}(A_2)=0{,}25$.\\
		Gọi $B$ là biến cố ``Linh kiện được chọn bị lỗi''.\\
		Theo giả thiết, tỉ lệ linh kiện bị lỗi của hai phân xưởng cho ta các xác suất có điều kiện: $\mathrm{P}(B|A_1)=0{,}04$ và $\mathrm{P}(B|A_2)=0{,}02$.
		\begin{itemchoice}
			\itemch Xác suất để linh kiện không bị lỗi biết rằng nó do phân xưởng $I$ sản xuất là
			\[\mathrm{P}(\overline{B}|A_1)=1-\mathrm{P}(B|A_1)=1-0{,}04=0{,}96.\]
			\itemch Xác suất để linh kiện được chọn bị lỗi và do phân xưởng $I$ sản xuất là
			\[\mathrm{P}(A_1B)=\mathrm{P}(A_1)\cdot\mathrm{P}(B|A_1)=0{,}75\cdot0{,}04=0{,}03.\]
			\itemch Áp dụng công thức xác suất toàn phần, xác suất để linh kiện được chọn bị lỗi là
			\[\mathrm{P}(B)=\mathrm{P}(A_1)\cdot\mathrm{P}(B|A_1)+\mathrm{P}(A_2)\cdot\mathrm{P}(B|A_2)=0{,}75\cdot0{,}04+0{,}25\cdot0{,}02=0{,}035.\]
			\itemch Áp dụng công thức Bayes, nếu linh kiện được chọn bị lỗi thì xác suất để nó do phân xưởng $I$ sản xuất là
			\[\mathrm{P}(A_1|B)=\dfrac{\mathrm{P}(A_1B)}{\mathrm{P}(B)}=\dfrac{0{,}03}{0{,}035}=\dfrac{30}{35}=\dfrac{6}{7}.\]
		\end{itemchoice}
	}
\end{ex}
\begin{ex}%[Thi thử SGD An Giang 25-26,L1]%[Lê Quốc Hiệp, DeTHPT2025-2026 ]%[2H5V4-4]
	Trong không gian $Oxyz$, cho mặt cầu $(S)$ có phương trình $(x-1)^2+(y-2)^2+(z-5)^2=25$ và đường thẳng $d$ đi qua hai điểm $A(7;8;5)$, $B(8;7;5)$. Giả sử mặt cầu $(S)$ được tịnh tiến đi lên (theo hướng dương của trục $Oz$) với phương vuông góc với mặt phẳng $(Oxy)$.
	\choiceTF
	{Bán kính của mặt cầu $(S)$ bằng $25$}
	{\True Vectơ $\vec{u}=(1;-1;0)$ là một vectơ chỉ phương của đường thẳng $d$}
	{\True Khoảng cách từ tâm mặt cầu $(S)$ ở vị trí ban đầu đến đường thẳng $d$ bằng $6\sqrt{2}$}
	{Có một thời điểm trong quá trình tịnh tiến, mặt cầu $(S)$ tiếp xúc với đường thẳng $d$}
	\loigiai{
		\begin{itemchoice}
			\itemch Mặt cầu $(S)$ có phương trình $(x-1)^2+(y-2)^2+(z-5)^2=25$ nên có tâm $I(1;2;5)$ và bán kính $R=\sqrt{25}=5\neq25$.
			\itemch Đường thẳng $d$ đi qua hai điểm $A(7;8;5)$ và $B(8;7;5)$ nên nhận vectơ $\overrightarrow{AB}=(8-7;7-8;5-5)=(1;-1;0)$ làm một vectơ chỉ phương.\\
			Do đó, $\vec{u}=(1;-1;0)$ là một vectơ chỉ phương của $d$.
			\itemch Tâm $I(1;2;5)$ ở vị trí ban đầu. Ta có $\overrightarrow{IA}=(6;6;0)$ và vectơ chỉ phương của $d$ là $\vec{u}=(1;-1;0)$.\\
			Tích có hướng $\left[\overrightarrow{IA},\vec{u}\right]=(0;0;-12)$.\\
			Khoảng cách từ $I$ đến $d$ là
			\[h=\dfrac{\left|\left[\overrightarrow{IA},\vec{u}\right]\right|}{|\vec{u}|}=\dfrac{\sqrt{0^2+0^2+(-12)^2}}{\sqrt{1^2+(-1)^2+0^2}}=\dfrac{12}{\sqrt{2}}=6\sqrt{2}.\]
			\itemch Giả sử mặt cầu $(S)$ tịnh tiến đi lên theo hướng dương của trục $Oz$ một đoạn $t>0$, khi đó tâm mặt cầu mới là $I_t(1;2;5+t)$ và bán kính không đổi $R=5$.\\
			Ta có $\overrightarrow{I_tA}=(6;6;-t)$. Tích có hướng $\left[\overrightarrow{I_tA},\vec{u}\right]=(-t;-t;-12)$.\\
			Khoảng cách từ $I_t$ đến đường thẳng $d$ là
			\begin{eqnarray*}
				d_t&=&\dfrac{\left|\left[\overrightarrow{I_tA},\vec{u}\right]\right|}{|\vec{u}|}\\
				&=&\dfrac{\sqrt{(-t)^2+(-t)^2+(-12)^2}}{\sqrt{1^2+(-1)^2+0^2}}\\
				&=&\dfrac{\sqrt{2t^2+144}}{\sqrt{2}}=\sqrt{t^2+72}.
			\end{eqnarray*}
			Mặt cầu tiếp xúc với đường thẳng $d$ khi và chỉ khi $d_t=R$, tức là
			\[\sqrt{t^2+72}=5\Leftrightarrow t^2+72=25\Leftrightarrow t^2=-47\text{ (vô nghiệm)}.\]
			Vậy không có thời điểm nào mặt cầu tiếp xúc với đường thẳng $d$.
		\end{itemchoice}
	}
\end{ex}
\Closesolutionfile{ans}

\caukq
\Opensolutionfile{ans}[ans/ans-56-26-2-AnGiang-SoGD-L1-KQ]
\begin{ex}%[Thi thử SGD An Giang 25-26,L1]%[Lê Quốc Hiệp, DeTHPT2025-2026 ]%[1H8V7-6]
	\immini{
		Cho khối chóp cụt tứ giác đều $ABCD.A'B'C'D'$ như hình vẽ. Biết tổng diện tích của hai mặt đáy bằng $54$ và độ dài đường chéo $AC'=9$. Tìm thể tích lớn nhất của khối chóp cụt đã cho.
	}{
		\usetikzlibrary{calc}
		
		\begin{tikzpicture}[font=\footnotesize,line join=round,line cap=round,]
			\def\k{1.3}
			% Đáy nhỏ
			\coordinate (A') at (0,0);
			\coordinate (B') at (-0.8,-1);
			\coordinate (D') at (1.8,0);
			\coordinate (C') at ($ (B') + (D') - (A') $);
			\coordinate (O') at ($ (B') + (D') - (A') $);
			% Đáy lớn
			\coordinate (A) at (105:2);
			\coordinate (B) at ($ (A) + \k*(B') - \k*(A') $);
			\coordinate (D) at ($ (A) + \k*(D') - \k*(A') $);
			\coordinate (C) at ($ (B) + (D) - (A) $);
			\draw[dashed]
			(A) node[above left]{$A$}
			-- (A') node[below]{$A'$}
			-- (B') node[left]{$B'$};
			\draw[dashed]
			(A') -- (D') node[right]{$D'$};
			\draw
			(A) -- (B) node[left]{$B$}
			-- (C) node[right]{$C$}
			-- (D) node[above right]{$D$}
			-- cycle;
			\draw
			(B) -- (B')
			-- (C') node[below]{$C'$}
			-- (D') -- (D);
			\draw (C) -- (C');
			\foreach \p in {A,B,C,D,A',B',C',D'}{
				\fill (\p) circle (1pt);
			}
		\end{tikzpicture}
	}
	\shortans[]{144}
	\loigiai{
		\immini
		{Gọi $x$, $y$ lần lượt là độ dài cạnh của hai đáy $ABCD$ và $A'B'C'D'$ ($x>0$, $y>0$).\\
		Theo giả thiết tổng diện tích hai mặt đáy bằng $54$, ta có $x^2+y^2=54$.\\
		Gọi $O$, $O'$ lần lượt là tâm của hai đáy $ABCD$ và $A'B'C'D'$. Chiều cao khối chóp cụt là $h=OO'$.\\
		Kẻ $C'H\perp AC$ ($H\in AC$), ta có $C'H=h$.\\
		Xét mặt cắt chéo $ACC'A'$ của khối chóp cụt, đây là một hình thang cân.\\
		Độ dài đoạn $AH$ được tính bằng nửa tổng hai đáy của hình thang cân.\\
		Đoạn $AH=\dfrac{AC+A'C'}{2}=\dfrac{x\sqrt{2}+y\sqrt{2}}{2}=\dfrac{\sqrt{2}(x+y)}{2}$.\\
		Trong tam giác vuông $AHC'$, ta có $AC'^2=AH^2+C'H^2$, suy ra
		\[h^2=9^2-\left(\dfrac{\sqrt{2}(x+y)}{2}\right)^2=81-\dfrac{(x+y)^2}{2}.\]}
		{\begin{tikzpicture}[font=\footnotesize,line join=round,line cap=round,]
				\def\k{1.3}
				% Đáy nhỏ
				\coordinate (A') at (0,0);
				\coordinate (B') at (-120:1);
				\coordinate (D') at (2,0);
				\coordinate (C') at ($ (B') + (D') - (A') $);
				\coordinate (O') at ($0.5*(C')$);
				% Đáy lớn
				\coordinate (A) at (105:2);
				\coordinate (B) at ($ (A) + \k*(B') - \k*(A') $);
				\coordinate (D) at ($ (A) + \k*(D') - \k*(A') $);
				\coordinate (C) at ($ (B) + (D) - (A) $);
				\coordinate (O) at ($(B)!0.5!(D)$);
				\coordinate (H) at ($ (O) + (C') - (O') $);
				\draw[dashed]
				(A) node[above left]{$A$}
				-- (A') node[left]{$A'$}
				-- (B') node[left]{$B'$};
				\draw[dashed]
				(C')--(A') -- (D') node[right]{$D'$} (A)--(C')--(H)node[above]{$H$};
				\draw[dashed]
				(O') node[below]{$O'$} -- (O) node[above]{$O$} ;
				\draw
				(A) -- (B) node[left]{$B$}
				-- (C) node[right]{$C$}
				-- (D) node[above right]{$D$}
				-- cycle;
				\draw
				(B) -- (B')
				-- (C') node[below]{$C'$}
				-- (D') -- (D);
				\draw (A)--(C) -- (C');
				\foreach \p in {A,B,C,D,A',B',C',D',O',O,H}{
					\fill (\p) circle (1pt);
				}
		\end{tikzpicture}}
		\noindent Đặt $t=xy$. Vì $(x-y)^2\ge0\Leftrightarrow2xy\le x^2+y^2\Rightarrow2t\le54\Leftrightarrow t\le27$.\\
		Do $x$, $y>0$ nên $t>0$. Vậy $t\in(0;27]$.\\
		Ta có $(x+y)^2=x^2+y^2+2xy=54+2t$.\\
		Suy ra $h^2=81-\dfrac{54+2t}{2}=54-t\Rightarrow h=\sqrt{54-t}$.\\
		Thể tích khối chóp cụt là
		\[V=\dfrac{1}{3}h(S_1+S_2+\sqrt{S_1S_2})=\dfrac{1}{3}\sqrt{54-t}(x^2+y^2+xy)=\dfrac{1}{3}\sqrt{54-t}(54+t).\]
		Xét hàm số $f(t)=(54-t)(54+t)^2$ với $t\in(0;27]$.
		\begin{eqnarray*}
			f'(t)&=&-(54+t)^2+2(54+t)(54-t)\\
			&=&(54+t)(-54-t+108-2t)\\
			&=&(54+t)(54-3t).
		\end{eqnarray*}
		Cho $f'(t)=0\Leftrightarrow54-3t=0\Leftrightarrow t=18$ (nhận).\\
		Lập bảng biến thiên, ta thấy hàm số $f(t)$ đạt giá trị lớn nhất tại $t=18$.\\
		Khi đó, thể tích lớn nhất của khối chóp cụt là
		\[V_{\max}=\dfrac{1}{3}\sqrt{54-18}(54+18)=\dfrac{1}{3}\sqrt{36}\cdot72=144.\]
	}
\end{ex}
\begin{ex}%[Thi thử SGD An Giang 25-26,L1]%[Lê Quốc Hiệp, DeTHPT2025-2026 ]%[2D4V3-4]
	\immini{
		Cho hàm số $f(x)=\dfrac{1}{4}x^3+x$. Đường thẳng $d$ đi qua gốc tọa độ $O$ và điểm $P(2;f(2))$ cắt đường thẳng $\Delta:y=-x+3$ tại điểm $Q$. Gọi $R$ là giao điểm của đường thẳng $\Delta$ với trục hoành. Gọi $A$ là diện tích hình phẳng giới hạn bởi đường cong $y=f(x)$, đường thẳng $\Delta$ và đoạn thẳng $PQ$; $B$ là diện tích hình phẳng giới hạn bởi đường cong $y=f(x)$, đường thẳng $\Delta$ và đoạn thẳng $OR$. Giá trị của $B-A$ là
	}{
		\begin{tikzpicture}[scale=1.25,font=\footnotesize, line join=round, line cap=round, >=stealth]
			\def\xS{1.25369}
			\def\yS{1.74631}
			% Trục tọa độ
			\draw[->] (-0.5,0) -- (3.7,0) node[right]{$x$};
			\draw[->] (0,-0.4) -- (0,4.7) node[above]{$y$};
			% Các điểm
			\coordinate (O) at (0,0);
			\coordinate (P) at (2,4);
			\coordinate (Q) at (1,2);
			\coordinate (R) at (3,0);
			\coordinate (S) at (\xS,\yS);
			% Miền B
			\fill[pattern=grid,pattern color=black!70](O) --plot[domain=0:\xS, samples=100](\x,{0.25*(\x)^3+\x})-- (R) -- cycle;
			% Miền A
			\fill[pattern=north west lines,pattern color=black!70](Q) -- (P) --plot[domain=2:\xS, samples=100](\x,{0.25*(\x)^3+\x})-- cycle;
			% Đồ thị y=f(x)
			\draw
			plot[domain=0:2.15, samples=120](\x,{0.25*(\x)^3+\x})node[above left]{$y=f(x)$};
			% Đường thẳng d: y=2x
			\draw (O) -- (2.25,4.5) node[right]{$d$};
			% Đường thẳng Delta: y=-x+3
			\draw (-0.1,3.1) -- (3.35,-0.35) node[below]{$\Delta$};
			% Nhãn miền
			\node[fill=white,inner sep=1pt] at (60:2.5) {$A$};
			\node[fill=white,inner sep=1pt] at (1.25,0.8) {$B$};
			% Điểm và nhãn
			\foreach \p/\pos in {O/below left,P/above left,Q/left,R/above right}{
				\fill (\p) circle (1pt) node[\pos]{$\p$};
			}
		\end{tikzpicture}
	}
	\shortans[]{2}
	\loigiai{
		Ta có $f(2)=\dfrac{1}{4}\cdot2^3+2=4\Rightarrow P(2;4)$.\\
		Đường thẳng $d$ đi qua gốc tọa độ $O(0;0)$ và $P(2;4)$ có phương trình là $y=2x$.\\
		Giao điểm $Q$ của $d$ và $\Delta$ thỏa mãn phương trình hoành độ giao điểm
		\[2x=-x+3\Leftrightarrow3x=3\Leftrightarrow x=1\Rightarrow y=2\Rightarrow Q(1;2).\]
		Giao điểm $R$ của $\Delta$ và trục hoành $Ox$ ($y=0$) thỏa mãn
		\[-x+3=0\Leftrightarrow x=3\Rightarrow R(3;0).\]
		Gọi $x_0$ là hoành độ giao điểm của đường cong $y=f(x)$ và đường thẳng $\Delta$. Dựa vào đồ thị, ta thấy $x_0\in(1;2)$.\\
		Diện tích $A$ là hình phẳng giới hạn bởi đường thẳng $d$, đường cong $y=f(x)$ và đường thẳng $\Delta$. Ta có
		\[A=\displaystyle\int\limits_1^{x_0}(2x-(-x+3))\,\mathrm{d}x+\displaystyle\int\limits_{x_0}^2(2x-f(x))\,\mathrm{d}x.\]
		Diện tích $B$ là hình phẳng giới hạn bởi đường cong $y=f(x)$, đường thẳng $\Delta$ và trục hoành $Ox$. Ta có
		\[B=\displaystyle\int\limits_0^{x_0}f(x)\,\mathrm{d}x+\displaystyle\int\limits_{x_0}^3(-x+3)\,\mathrm{d}x.\]
		Xét hiệu $B-A$:
		\begin{eqnarray*}
			B-A&=&\displaystyle\int\limits_0^{x_0}f(x)\,\mathrm{d}x+\displaystyle\int\limits_{x_0}^3(-x+3)\,\mathrm{d}x-\displaystyle\int\limits_1^{x_0}(3x-3)\,\mathrm{d}x-\displaystyle\int\limits_{x_0}^2(2x-f(x))\,\mathrm{d}x\\
			&=&\displaystyle\int\limits_0^{x_0}f(x)\,\mathrm{d}x+\displaystyle\int\limits_{x_0}^3(-x+3)\,\mathrm{d}x-\left(\displaystyle\int\limits_1^{x_0}2x\,\mathrm{d}x-\displaystyle\int\limits_1^{x_0}(-x+3)\,\mathrm{d}x\right)-\displaystyle\int\limits_{x_0}^22x\,\mathrm{d}x+\displaystyle\int\limits_{x_0}^2f(x)\,\mathrm{d}x\\
			&=&\left(\displaystyle\int\limits_0^{x_0}f(x)\,\mathrm{d}x+\displaystyle\int\limits_{x_0}^2f(x)\,\mathrm{d}x\right)+\left(\displaystyle\int\limits_1^{x_0}(-x+3)\,\mathrm{d}x+\displaystyle\int\limits_{x_0}^3(-x+3)\,\mathrm{d}x\right)-\left(\displaystyle\int\limits_1^{x_0}2x\,\mathrm{d}x+\displaystyle\int\limits_{x_0}^22x\,\mathrm{d}x\right)\\
			&=&\displaystyle\int\limits_0^2f(x)\,\mathrm{d}x+\displaystyle\int\limits_1^3(-x+3)\,\mathrm{d}x-\displaystyle\int\limits_1^22x\,\mathrm{d}x.
		\end{eqnarray*}
		Vậy $B-A=3+2-3=2$.
	}
\end{ex}
\begin{ex}%[Thi thử SGD An Giang 25-26,L1]%[Lê Quốc Hiệp, DeTHPT2025-2026 ]%[0D0V2-7]
	Trong mặt phẳng tọa độ $Oxy$, xét tập hợp $S$ gồm các điểm có tọa độ $(x;y)$ với $x$, $y$ là các số nguyên dương không vượt quá $15$. Chọn ngẫu nhiên $2$ điểm phân biệt $A$, $B$ từ tập $S$. Gọi $C$ là trung điểm của đoạn thẳng $AB$. Gọi $p$ là xác suất để điểm $C$ có tọa độ đều là các số nguyên. Tính giá trị của $225p$.
	\shortans[]{56}
	\loigiai{
		Từ giả thiết, $x$, $y\in\{1;2;\ldots;15\}$.\\
		Tập $S$ có tổng cộng $15\cdot15=225$ điểm.\\
		Không gian mẫu là chọn ngẫu nhiên $2$ điểm phân biệt từ $225$ điểm của $S$.\\
		Số phần tử của không gian mẫu là $n(\Omega)=C_{225}^2=25200$.\\
		Gọi biến cố $E$: ``Điểm $C$ có tọa độ đều là các số nguyên''.\\
		Vì $C$ là trung điểm của $AB$ nên $x_C=\dfrac{x_A+x_B}{2}$ và $y_C=\dfrac{y_A+y_B}{2}$.\\
		Để $x_C$, $y_C\in\mathbb{Z}$ thì $x_A+x_B$ và $y_A+y_B$ phải là các số chẵn.\\
		Điều này tương đương với $x_A$, $x_B$ cùng tính chẵn lẻ và $y_A$, $y_B$ cùng tính chẵn lẻ.\\
		Trong tập $\{1;2;\ldots;15\}$ có $8$ số lẻ và $7$ số chẵn.\\
		Ta phân chia $225$ điểm của tập $S$ thành $4$ nhóm theo tính chẵn lẻ của hoành độ và tung độ:
		\begin{itemize}
			\item Nhóm 1 (lẻ; lẻ): Có $8\cdot8=64$ điểm.
			\item Nhóm 2 (chẵn; chẵn): Có $7\cdot7=49$ điểm.
			\item Nhóm 3 (lẻ; chẵn): Có $8\cdot7=56$ điểm.
			\item Nhóm 4 (chẵn; lẻ): Có $7\cdot8=56$ điểm.
		\end{itemize}
		Hai điểm $A$, $B$ thỏa mãn điều kiện khi và chỉ khi chúng được chọn cùng thuộc một trong bốn nhóm trên.\\
		Số kết quả thuận lợi cho biến cố $E$ là
		\[n(E)=C_{64}^2+C_{49}^2+C_{56}^2+C_{56}^2=6272.\]
		Xác suất cần tìm là
		\[p=\dfrac{n(E)}{n(\Omega)}=\dfrac{6272}{25200}=\dfrac{56}{225}.\]
		Vậy giá trị của $225p$ là
		\[225p=225\cdot\dfrac{56}{225}=56.\]
	}
\end{ex}
\begin{ex}%[Thi thử SGD An Giang 25-26,L1]%[Lê Quốc Hiệp, DeTHPT2025-2026 ]%[2H5V4-4]
	Trong không gian với hệ trục tọa độ $Oxyz$, đơn vị trên mỗi trục tọa độ là mét, một thiết bị phát tia laser được đặt tại điểm $A(5\sqrt{3};8;15)$. Thiết bị này chiếu một tia sáng về phía bức tường phẳng trùng với mặt phẳng tọa độ $(Oyz)$. Biết rằng tia sáng phát ra luôn thay đổi nhưng luôn tạo với trục $Ox$ một góc $60^\circ$. Gọi $M$ là vị trí vệt sáng laser chiếu lên bức tường. Khoảng cách lớn nhất từ vệt sáng $M$ đến gốc tọa độ $O$ bằng bao nhiêu mét?
	\shortans[]{32}
	\loigiai{
		\immini
		{Vệt sáng $M$ nằm trên bức tường trùng với mặt phẳng $(Oyz)$ nên tọa độ của $M$ có dạng $M(0;y;z)$.\\
		Trục $Ox$ có một vectơ chỉ phương là $\vec{i}=(1;0;0)$.\\
		Ta có $\overrightarrow{AM}=(-5\sqrt{3};y-8;z-15)$.\\
		Vì tia sáng luôn tạo với trục $Ox$ một góc $60^\circ$ nên ta có
		\begin{eqnarray*}
			&&\left|\cos(\overrightarrow{AM},\vec{i})\right|=\cos60^\circ=\dfrac{1}{2}\\
			&\Leftrightarrow&\dfrac{\left|\overrightarrow{AM}\cdot\vec{i}|\right|}{\left|\overrightarrow{AM}\right|\cdot\left|\vec{i}\right|}=\dfrac{1}{2}\\
			&\Leftrightarrow&\dfrac{\left|-5\sqrt{3}\cdot1+0+0\right|}{\sqrt{(-5\sqrt{3})^2+(y-8)^2+(z-15)^2}\cdot1}=\dfrac{1}{2}\\
			&\Leftrightarrow&\dfrac{5\sqrt{3}}{\sqrt{75+(y-8)^2+(z-15)^2}}=\dfrac{1}{2}\\
			&\Leftrightarrow&\sqrt{75+(y-8)^2+(z-15)^2}=10\sqrt{3}\\
			&\Leftrightarrow&75+(y-8)^2+(z-15)^2=300\\
			&\Leftrightarrow&(y-8)^2+(z-15)^2=225.
		\end{eqnarray*}}
		{\begin{tikzpicture}[font=\footnotesize, line join=round, line cap=round, >=stealth, color=blue]
				% Các trục tọa độ
				\draw[->] (0,0) -- (4,0) node[below] {$x$};
				\draw[->] (0,0) -- (-2.5,-2) node[left] {$y$};
				\draw[->] (0,0) -- (0,3.5) node[right] {$z$};
				\node[below right] at (0,0) {$O$};
				
				% Định nghĩa các điểm
				\coordinate (H) at (-1.5, 1.5);
				\coordinate (A) at (2, 1.5);
				\coordinate (M) at (-1.5, 3);
				\coordinate (M') at (-1.5, 0);
				\coordinate (O) at (0, 0);
				% Vẽ đường tròn đáy hình nón (Elip trong phối cảnh)
				\draw (-1.5, 3) arc (90:270:0.5 and 1.5);
				\draw[dashed] (-1.5, 3) arc (90:-90:0.5 and 1.5);
				\draw (-1.5, 3) arc (90:270:0.5 and 1.5)
				coordinate[pos=0.5] (P);
				% Các đường đứt nét bên trong khối
				\draw[dashed] (H) -- (M);
				\draw[dashed] (H) -- (M');
				\draw[dashed] (H) -- (A);
				
				% Ký hiệu góc vuông tại H
				\draw ($ (H) + (0.15, 0) $) -- ($ (H) + (0.15, 0.15) $) -- ($ (H) + (0, 0.15) $);
				
				% Vẽ và tô màu góc 60 độ
				% Tính góc của đường AM so với phương ngang
				\pgfmathsetmacro{\angleAM}{180-atan(1.5/3.5)} 
				\fill[blue!20] (A) -- ++(-0.8,0) arc (180:\angleAM:0.8) -- cycle;
				\draw (A) ++(-0.8,0) arc (180:\angleAM:0.8);
				\node at ($ (A) + (-1.1, 0.25) $) {$60^\circ$};
				
				% Các đường sinh và đường thẳng mở rộng
				\draw (M) -- (A) -- ($ (A) ! -0.4 ! (M) $);
				\draw (M') -- (A) -- ($ (A) ! -0.5 ! (M') $) node[above left] {$\Delta$};
				\draw (A) -- ($ (A) + (1.5, 0) $) node[below] {$d$};
				\path[name path=OP] (O) -- (P);
				\path[name path=AM] (A) -- (M');
				% Tìm giao điểm
				\path[name intersections={of=OP and AM, by=I}];
				% Vẽ điểm giao
				\node[shift={(120:2ex)}] at (H){$H$};;
				\node[below] at (A) {$A$};
				\node[left] at (P) {$M$};
				\draw[dashed] (I)--(P);
				\draw (I)--(O);
				\foreach \p in {O,H,A,P,I}
				{
					\fill (\p) circle (1pt);
				}
		\end{tikzpicture}}
		\noindent Suy ra tập hợp các điểm $M$ là một đường tròn $(C)$ nằm trong mặt phẳng $(Oyz)$, có tâm $I(0;8;15)$ và bán kính $R=\sqrt{225}=15$.\\
		Gốc tọa độ $O(0;0;0)$ nằm trên mặt phẳng $(Oyz)$.\\
		Khoảng cách từ $O$ đến tâm $I$ là
		\[OI=\sqrt{0^2+8^2+15^2}=17.\]
		Vì $M$ thuộc đường tròn $(C)$ tâm $I$, bán kính $R$ nên khoảng cách lớn nhất từ $M$ đến $O$ là
		\[OM_{\max}=OI+R=17+15=32.\]
		Vậy khoảng cách lớn nhất từ vệt sáng $M$ đến gốc tọa độ $O$ bằng $32$ mét.
	}
\end{ex}
\begin{ex}%[Thi thử SGD An Giang 25-26,L1]%[Lê Quốc Hiệp, DeTHPT2025-2026 ]%[1D6V4-6]
	Một nhóm phát triển phần mềm độc lập cung cấp một ứng dụng Trí tuệ nhân tạo hỗ trợ quản lý công việc và học tập cá nhân. Mỗi người dùng đăng ký bản quyền sử dụng sẽ phải trả mức phí cố định là $2$ triệu đồng/năm. Tổng chi phí duy trì hệ thống lưu trữ đám mây và cập nhật phần mềm trong một năm phụ thuộc vào số lượng người dùng $x$ ($x\in\mathbb{N}^*$) đang sử dụng và được mô hình hóa bởi hàm số $C(x)=10\ln(x)+50$ (triệu đồng). Để nhóm phát triển đạt mức lợi nhuận tối thiểu là $200$ triệu đồng trong một năm, họ cần thu hút được ít nhất bao nhiêu người dùng đăng ký phần mềm đó?
	\shortans[]{151}
	\loigiai{
		Gọi $x$ ($x\in\mathbb{N}^*$) là số lượng người dùng đăng ký.\\
		Tổng doanh thu trong một năm của nhóm là $R(x)=2x$ (triệu đồng).\\
		Lợi nhuận trong một năm nhóm phát triển thu được là
		\[P(x)=R(x)-C(x)=2x-(10\ln(x)+50)=2x-10\ln(x)-50.\]
		Để đạt mức lợi nhuận tối thiểu $200$ triệu đồng, ta phải có
		\[P(x)\ge200\Leftrightarrow2x-10\ln(x)-50\ge200\Leftrightarrow2x-10\ln(x)-250\ge0.\]
		Xét hàm số $f(x)=2x-10\ln(x)-250$ trên khoảng $(0;+\infty)$.\\
		Ta có $f'(x)=2-\dfrac{10}{x}=\dfrac{2x-10}{x}$.\\
		Với $x\ge5$ thì $f'(x)\ge0$, do đó hàm số $f(x)$ đồng biến trên nửa khoảng $[5;+\infty)$.\\
		Vì $x\in\mathbb{N}^*$ nên ta có thể kiểm tra các giá trị nguyên của $x$:
		\begin{itemize}
			\item Với $x=150$, ta có $f(150)=2\cdot150-10\ln(150)-250=50-10\ln(150)\approx-0{,}106<0$ (không thỏa mãn).
			\item Với $x=151$, ta có $f(151)=2\cdot151-10\ln(151)-250=52-10\ln(151)\approx1{,}827>0$ (thỏa mãn).
		\end{itemize}
		Do hàm số đồng biến khi $x\ge5$ nên $f(x)\ge0$ với mọi số nguyên $x\ge151$.\\
		Vậy họ cần thu hút ít nhất $151$ người dùng đăng ký phần mềm.
	}
\end{ex}
\begin{ex}%[Thi thử SGD An Giang 25-26,L1]%[Lê Quốc Hiệp, DeTHPT2025-2026 ]%[1D2V3-8]
	Đầu năm $2025$, ông A thành lập một doanh nghiệp. Tổng số tiền ông A dùng để trả chi phí vận hành trong năm $2025$ là $2$ tỷ đồng. Biết rằng cứ sau mỗi năm thì tổng số tiền dùng để trả chi phí vận hành trong năm đó tăng thêm $10\%$ so với năm trước. Hỏi năm nào là năm đầu tiên mà tổng số tiền ông A dùng để trả chi phí vận hành trong cả năm lớn hơn $4$ tỷ đồng?
	\shortans[]{2033}
	\loigiai{
		Gọi $C_n$ là tổng số tiền ông A dùng để trả chi phí vận hành trong năm thứ $n$ kể từ năm $2025$ (với $n\in\mathbb{N}^*$, $n=1$ ứng với năm $2025$).\\
		Theo giả thiết, số tiền chi phí vận hành mỗi năm tạo thành một cấp số nhân với số hạng đầu $C_1=2$ và công bội $q=1+10\%=1{,}1$.\\
		Số tiền chi phí vận hành trong năm thứ $n$ là
		\[C_n=C_1\cdot q^{n-1}=2\cdot1{,}1^{n-1}.\]
		Để tổng số tiền này lớn hơn $4$ tỷ đồng, ta có bất phương trình
		\begin{eqnarray*}
			&&2\cdot1{,}1^{n-1}>4\\
			&\Leftrightarrow&1{,}1^{n-1}>2\\
			&\Leftrightarrow& n-1>\log_{1{,}1}2\\
			&\Leftrightarrow& n>1+\log_{1{,}1}2\approx8{,}27.
		\end{eqnarray*}
		Vì $n$ là số nguyên dương nên giá trị nhỏ nhất của $n$ thỏa mãn là $n=9$.\\
		Năm đầu tiên tương ứng là $2025+(9-1)=2033$.\\
		Vậy năm $2033$ là năm đầu tiên chi phí vận hành lớn hơn $4$ tỷ đồng.
	}
\end{ex}
\Closesolutionfile{ans}

\indapan{6}{ans/ans-56-26-2-AnGiang-SoGD-L1-LC}
\indapan{4}{ans/ans-56-26-2-AnGiang-SoGD-L1-DS}
\indapan{6}{ans/ans-56-26-2-AnGiang-SoGD-L1-KQ}
